
This is an interactive problem. Your submission will be run against an interactor, which reads
the standard output of your submission and writes to the standard input of your submission. This
interaction needs to follow a specific protocol:

The interactor first sends
one line with an integer $n$ ($1 \leq n \leq 2\cdot 10^5$),
indicating that the domain of the hidden function $f$ is $\{1,\dots,n\}$.

Then, your program should make at most $1000$ queries to find the answer.
Each query is made by printing one line of the form ``\texttt{?~$c$~$r$}'' ($1 \leq c, r \leq n$).
The interactor will respond with an integer $h$ ($1 \leq h \leq n$), the value of $f^c(r)$.

When you have determined some values of $c$ and $r$ such that $f^c(r) = c$,
print one line of the form ``\texttt{!~$c$~$r$}'' ($1 \leq c, r \leq n$),
after which the interaction will stop.
Printing the answer does not count as a query.

If there are multiple valid solutions, you may output any one of them.

The interactor is not adaptive: the hidden function $f$ is fixed up front, and does not depend on your queries.

Make sure you flush the buffer after each write.

A testing tool is provided to help you develop your solution.

Using more than $1000$ queries will result in a wrong answer.
